\documentclass[11pt,a4paper]{article}
\usepackage[T1]{fontenc}
\usepackage{lmodern}
\usepackage[margin=25mm]{geometry}
\usepackage{amsmath,amssymb,amsthm,booktabs,array,microtype}
\usepackage[hidelinks]{hyperref}
\hypersetup{pdftitle={An exact mean bound for the rational complexity of binary strings},pdfauthor={Oleksiy Babanskyy}}
\numberwithin{equation}{section}
\newtheorem{theorem}{Theorem}[section]
\newtheorem{lemma}[theorem]{Lemma}
\newtheorem{proposition}[theorem]{Proposition}
\theoremstyle{remark}
\newtheorem{remark}[theorem]{Remark}
\newcommand{\E}{\mathbb E}
\newcommand{\Z}{\mathbb Z}
\newcommand{\pos}[1]{\left[#1\right]_+}
\newcommand{\odd}{\mathrm{odd}}
\title{An exact mean bound for the rational complexity of binary strings}
\author{Oleksiy Babanskyy}
\date{}
\begin{document}
\maketitle
\begin{abstract}
For a residue $r$ modulo $2^N$, let $R_N(r)$ be the least height
$\max(|a|,b)$ of a rational representation $a\equiv rb\pmod{2^N}$
with positive odd denominator $b$.
We prove that the uniform mean of the ordinary, unrounded quantity
$\log_2R_N(r)$ is at most $N/2$ for every $N\ge1$.
This is the mean part of Question~3.1 in the 2026 survey of Lin, Xiao and Chen.
The proof compares the product of the two lift heights with a truncated
logarithmic excess at the parent modulus. Short lattice vectors with even
denominator cover the large-height residues; completion in the actual odd
coset bounds their contribution. A finite lattice sum, a two-prime sieve
and an explicit primitive-point estimate close the inequality for parent
length at least $26$. Integer certificates cover the remaining parent
lengths $1,\ldots,25$. The variance question is left open.
\end{abstract}

\section{The problem and the result}
Write a binary string as $s=(s_0,\ldots,s_{N-1})$, and identify it with
$r=\sum_{i=0}^{N-1}s_i2^i$ modulo $2^N$. Define
\begin{equation}\label{eq:def}
 R_N(r)=\min_{\substack{a\equiv rb\pmod{2^N}\\b>0,\ b\ \odd}}
                   \max(|a|,b),\qquad
 \mathcal E_N=2^{-N}\sum_{r=0}^{2^N-1}\log_2R_N(r).
\end{equation}
This is the rational height underlying ordinary finite $2$-adic complexity.
The logarithm is not rounded to an integer. The measure includes every
string, including the zero string; it is not a measure on odd residues.

\begin{theorem}\label{thm:main}
For every integer $N\ge1$,
\[
                         \mathcal E_N\le \frac N2.
\]
\end{theorem}

Tian and Qi studied this same mean and proved a weaker upper bound
\cite[Theorem~3]{TQ}. Chen and Winterhof proved the lower bound
$\mathcal E_N\ge N/2-1$ and an upper bound with a logarithmic error
\cite[Theorems~2--3]{CWv1}. Their short primitive counts credit Tian--Qi;
their growth recurrence credits Goresky--Klapper \cite{GK}.
The exact upper in Theorem~\ref{thm:main} is explicitly asked for in
\cite[Question~3.1]{LXC}, together with a separate variance problem.
We answer only the mean part.

Our proof uses elementary geometry of congruence lattices and finite
integer certificates. It does not use a limiting distribution to infer
an inequality at every length. Its computational premises are stated in
Section~\ref{sec:finite} and checked by the accompanying dependency-free
program. The uniform argument is a proof for all sufficiently large
lengths, not an extrapolation of those certificates.

\paragraph{Source and priority boundary.}
For Chen--Winterhof's detailed statements we use the explicitly identified
arXiv version~1 \cite{CWv1}. The final publication \cite{CWfinal} has been
identified, but its full body has not been obtained for this paper. We do not assert a lemma-by-lemma identity of those versions.
The July 2026 survey, which includes Chen as an author, explicitly poses
the exact mean question after the final publication.
No exact antecedent has been identified in the primary sources checked
for this paper; this is not a certification of universal priority.
Ren--Winterhof's recent symmetric complexity \cite{RW} minimizes also
over a reversed string and is a different random variable.

\paragraph{Notation.}
All logarithms in the proof are natural. The subscript in $R_N$ denotes
a length. When a parent modulus is $M=2^n$, abbreviate $R_M(r)=R_n(r)$;
thus $R_{2M}$ denotes the child height $R_{n+1}$. Put
\[
 q=\sqrt M,\qquad s=q/\sqrt2,\qquad J=\lfloor n/2\rfloor.
\]
The convention $[x]_+=\max(x,0)$ will be used.
A minimum in \eqref{eq:def} is primitive in $\Z^2$: its common divisor
is odd and invertible modulo $2^N$, so division would otherwise lower
its height. In particular $R_M(0)=1$ and $R_M(r)\le M/2$ for $M\ge2$.

\section{The two lifts and a logarithmic budget}
\begin{lemma}\label{lem:children}
Let $M=2^n$, $n\ge1$. For a residue $r$ modulo $M$, put
\[
 h=R_M(r),\qquad H=\max\{R_{2M}(r),R_{2M}(r+M)\}.
\]
Then $\min\{R_{2M}(r),R_{2M}(r+M)\}=h$, and
\begin{equation}\label{eq:product}
                            hH\le\max(M,2h^2).
\end{equation}
\end{lemma}
\begin{proof}
The union of the two sets of child representations is exactly the set of
parent representations, since an odd denominator is invertible modulo
$2M$. Choose a primitive parent minimum $u=(a,b)$, $b>0$ odd.
It selects a child $r_0$ with height $h$.
Reflection allows $a\ge0$.

Suppose first $a>0$. Integer solutions $v=(c,d)$ of
$ad-bc=M$ form $v_0+ku$ by B\'ezout's identity. The restriction $d$ odd
selects exactly one class of $k$ modulo $2$. Moreover
\[
 c\equiv(r_0+M)d\pmod{2M}.
\]
Indeed $a\equiv r_0b\pmod{2M}$ and $b^{-1}M\equiv M\pmod{2M}$,
while $Md\equiv M$ for odd $d$. A negative $d$ can be oriented positive
by negating $v$; its height and child are unchanged.

Put $\tau=\min(a,b)$, $h=\max(a,b)$, $\sigma=a+b$ and
$X=\max\{h,(M+2ab)/\sigma\}$.
Parameterize the real determinant line by
$(-M/\sigma,M/\sigma)+t(a,b)$.
If $M\le h(h-\tau)$, then $X=h$. When $a=h$, the full interval
$c\in[-h,h]$ gives $|d|=|(M+bc)/a|\le h$; its length in $t$ is $2$.
When $b=h$, use the full interval $d\in[-h,h]$ instead.
If $M\ge h(h-\tau)$, then $X=(M+2ab)/\sigma\ge h$. The points at
$t=-2b/\sigma$ and $t=2a/\sigma$ are
\[
 \left(-X,\frac{M-2b^2}{\sigma}\right),\qquad
 \left(\frac{-M+2a^2}{\sigma},X\right).
\]
Both belong to $[-X,X]^2$: the required conditions are
$M\ge b(b-a)$ and $M\ge a(a-b)$, which follow from the case assumption.
The closed segment between them has parameter length $2$.
In either case it intersects the class $t_0+2\Z$.
Thus the other child has height at most $X$.

Writing $x=\tau/h$, we obtain
\[
 hH\le\max\left\{h^2,\frac{M+2h^2x}{1+x}\right\}
       \le\max(M,2h^2).
\]
If $a=0$, primitivity forces $b=1$, the children are $0$ and $M$, and
their heights are $1$ and $M$, so the same conclusion holds.
Ties between parent minima need no special choice.
\end{proof}

Define the exact finite quantities
\begin{equation}\label{eq:budgets}
\begin{split}
 C_M(t)&=\#\{r\bmod M:R_M(r)\le t\},\\
 K_M&=\frac1M\sum_r\pos{\ln(R_M(r)/s)},\\
 \mathcal D_M&=\frac1M\int_s^q C_M(t)\,\frac{dt}{t},\\
 \mathcal L_M&=\frac1M\sum_r\pos{\ln(R_M(r)/q)}.
\end{split}
\end{equation}
We reserve $D$ without a calligraphic letter for a dyadic content below.
Taking logarithms in \eqref{eq:product}, and averaging over both children,
gives
\begin{equation}\label{eq:reduction}
 \E_{r\bmod2M}\ln R_{2M}(r)\le \frac12\ln M+K_M.
\end{equation}
There is no induction hypothesis on the parent mean.

For every $R>0$,
\[
 \pos{\ln(R/s)}
 =\ln(q/s)-\int_s^q {\bf1}_{R\le t}\frac{dt}{t}
             +\pos{\ln(R/q)}.
\]
Consequently
\begin{equation}\label{eq:identity}
 K_M=\frac12\ln2-\mathcal D_M+\mathcal L_M.
\end{equation}
It suffices to prove $K_M\le\frac12\ln2$ at every parent length.
For large lengths we prove $\mathcal L_M<\mathcal D_M$.

\section{Short even vectors and the actual odd coset}\label{sec:cusps}
For $r\bmod M$ let
\[
 \Lambda_r=\{(A,B)\in\Z^2:A-rB\in M\Z\}.
\]
Its basis $(M,0),(r,1)$ has determinant $M$.

\begin{lemma}\label{lem:certificate}
If $R_M(r)>q$, there exist
\[
 D=2^j,\quad1\le j\le J,\quad b>0\ \odd,\quad
 \gcd(a,b)=1,\quad D\max(|a|,b)\le q
\]
such that
\begin{equation}\label{eq:contact}
           a-rb=(M/D)\alpha_r,\qquad \alpha_r\ \odd.
\end{equation}
For each such triple, its cell $\mathcal C(D,a,b)$ of residues satisfying
\eqref{eq:contact} has exactly $D/2$ elements.
For each residue in that cell, $D(a,b)$ is primitive as an element of
$\Lambda_r$.
\end{lemma}
\begin{proof}
Minkowski's convex body theorem \cite{Cassels} in the square of radius
$q+\varepsilon$ gives a nonzero lattice vector of height at most
$q+\varepsilon$. Discreteness permits passage to the closed radius $q$.
Choose a shortest nonzero vector. Its denominator is nonzero, since a
nonzero horizontal vector has height at least $M>q$.
Orient it positive. The hypothesis $R_M(r)>q$ forces its denominator
even, and its numerator is then even as well.
A common odd factor could be divided out while preserving lattice
membership, contradicting minimality. Thus its content is $D=2^j\ge2$;
write it as $D(a,b)$ with $\gcd(a,b)=1$.
The congruence modulo $Q=M/D$, which is even, forces $b$ odd.
If $\alpha_r=(a-rb)/Q$ were even, division by $2$ would again yield
a shorter vector of $\Lambda_r$. This proves the asserted contact.

There is one solution of $a-rb\equiv0\pmod Q$. Its $D$ lifts modulo $M$
change $\alpha_r$ by $-bt$, $t\bmod D$. Exactly $D/2$ have odd contact.
The coordinates of $D(a,b)$ in the displayed lattice basis are
$(\alpha_r,Db)$. Their gcd is $1$: $\alpha_r$ is odd, and
$\gcd(\alpha_r,b)=1$ follows from $\gcd(a,b)=1$ and $\gcd(Q,b)=1$.
This is lattice primitivity, despite the even integer content.
\end{proof}

\begin{lemma}\label{lem:completion}
On a cell $\mathcal C(D,a,b)$ put $A=D|a|$, $B=Db$. Then
\begin{equation}\label{eq:X}
 R_M(r)\le X:=\frac{M+AB}{A+B},\qquad X\ge q.
\end{equation}
On the axis $a=0,b=1$, the height is exactly $R_M(r)=M/D$.
\end{lemma}
\begin{proof}
For $a>0$, lattice primitivity allows completion of $v=(A,B)$ to a
basis with determinant $M$. Such completions $w=(c,d)$ satisfy
\[
 ad-bc=Q,\qquad ad\equiv Q\pmod b,\qquad
                       \alpha_r d\equiv1\pmod D.
\]
The last two congruences give one class of odd $d$ modulo $B=Db$ by
the Chinese remainder theorem. Invert $a$ only modulo $b$; no inversion
modulo an even denominator is needed. For $b=1$ the first congruence
is empty. All completions differ by a multiple of $v$.

Put $x=M/(A+B)$. Let $d_-\le x<d_+=d_-+B$ be the adjacent denominators
in this class, and write $\eta=x-d_-\in[0,B)$.
On the determinant line, $c(d)=(Ad-M)/B$. The two heights are
\begin{equation}\label{eq:adjacent}
 \max(|c(d_-)|,|d_-|)=x+(A/B)\eta,\qquad
 \max(|c(d_+)|,|d_+|)=x+B-\eta.
\end{equation}
For clarity, the sign checks use $A,B\le q$.
If $d_-<0$ and $A<B$, then
$(B-A)|d_-|\le(B-A)(B-x)\le B^2\le M$, so $-c(d_-)\ge|d_-|$;
the case $A\ge B$ is immediate.
At $d_+>x$, $-c(d_+)\le d_+$.
If $A>B$, then $c(d_+)\le d_+$ follows from
$(A-B)(x+B)\le M$, equivalent to $A^2-B^2\le2M$.
The other sign cases follow directly from $d_-\le x<d_+$.
The smaller expression in \eqref{eq:adjacent} is at most
$x+AB/(A+B)=X$, since the two affine expressions meet at
$\eta=B^2/(A+B)$.
Negative odd denominators are oriented positive by negating both
coordinates. Reflection handles $a<0$.
Finally $X/q\ge1$ is equivalent to
$(1-A/q)(1-B/q)\ge0$.

If $a=0$, primitivity gives $b=1$ and $r=-Q\alpha_r$.
Every marked numerator is a nonzero odd multiple of $Q$, so the height
is at least $Q$.
Take $c=-Q$ and a positive odd $d\le D-1$ representing
$\alpha_r^{-1}\pmod D$. This has height $Q$, because $D\le q\le Q$.
\end{proof}

\begin{remark}
Only coverage and the cell masses are needed for our upper bound.
We do not require disjointness of the cells or replace their actual
odd phases by continuous phase averages. The union bound below remains
valid even without a uniqueness statement.
\end{remark}

Let $V=\{r\ne0:v_2(r)\ge\lceil n/2\rceil\}$.
The axis cells are precisely these residues: for $D=2^j$ their valuation
is $k=n-j$ and their height is $2^k$.
This family is already present in Tian--Qi \cite[Lemma~4]{TQ}.
Write $\mathcal L_V$ for its contribution to $\mathcal L_M$ and
$\mathcal L_{\rm other}=\mathcal L_M-\mathcal L_V$.
For $0<x,y\le1$ define
\[
 f(x,y)=\ln\frac{1+xy}{x+y},\qquad
 P_f(T)=\sum_{\substack{1\le a,b\le T\\b\ \odd,\ \gcd(a,b)=1}}
                                           f(a/T,b/T).
\]
The kernel is nonnegative. Every non-axis residue with height above $q$
is covered by Lemma~\ref{lem:certificate}, and Lemma~\ref{lem:completion}
bounds its energy by $f(D|a|/q,Db/q)$.
Each cell has $D/2$ residues and both signs of $a$ give the weight $D$.
Therefore, using an upper bound by a union of cells,
\begin{equation}\label{eq:U}
 \mathcal L_{\rm other}\le U_M:=\frac1M
                         \sum_{j=1}^J2^jP_f(q/2^j).
\end{equation}
Including cells whose height is at most $q$ only enlarges this upper
bound; all summands are nonnegative.

\section{An explicit estimate of the lattice sum}
Let
\[
 I=\int_{[0,1]^2}f(x,y)\,dx\,dy,\qquad
 F(T)=\sum_{\substack{1\le a,b\le T\\b\ \odd}}f(a/T,b/T).
\]
The integral is finite: $0\le f(x,y)\le-\ln\max(x,y)$, so $I\le1/2$.

\begin{lemma}\label{lem:grid}
For every real $T>0$,
\begin{equation}\label{eq:grid}
                   \frac I2T^2-T\le F(T)\le\frac I2T^2,
 \qquad I=\frac{\zeta(2)-1}{2}=\frac{\pi^2}{12}-\frac12.
\end{equation}
\end{lemma}
\begin{proof}
For fixed $0<x\le1$, extend $f(x,y)$ to $y\ge0$. Its derivatives satisfy
\[
 \partial_yf=\frac{x^2-1}{(1+xy)(x+y)}\le0,\qquad
 \partial_y^2f=\frac1{(x+y)^2}-\frac{x^2}{(1+xy)^2}\ge0.
\]
The last assertion follows from $1+xy\ge x(x+y)$, also beyond $y=1$.
Odd $b\le T$ are the midpoints of intervals
$[(b-1)/T,(b+1)/T]$. Convexity gives
\[
 \frac2T\sum_{b\le T,\ b\ \odd} f(x,b/T)
 \le\int_0^{Y_T}f(x,y)\,dy,\qquad
 Y_T=\frac{2\lceil\lfloor T\rfloor/2\rceil}{T}.
\]
If $Y_T\le1$, the omitted portion is nonnegative.
If $Y_T>1$, the added portion is nonpositive, because $f(x,y)\le0$
for $y>1$. Thus this integral is at most
$\Phi(x)=\int_0^1f(x,y)\,dy$.
The function $\Phi$ is decreasing and nonnegative, so its right
Riemann sum at $a/T$, $1\le a\le\lfloor T\rfloor$, is at most $TI$.
This proves the upper bound for real $T$, including the last odd cell.

For $T\ge1$, the clipped rectangles
$[a/T,(a+1)/T]\times[b/T,(b+2)/T]$ with the same indices cover
$[1/T,1]^2$ with disjoint interiors. Monotonicity bounds the integral
there by $2F(T)/T^2$.
The omitted vertical and horizontal strips have total integral at most
$2/T$, using $f\le-\ln y$ and $f\le-\ln x$, respectively.
This proves the lower bound. For $0<T<1$, $F(T)=0$ and the lower bound
is nonpositive.

For $z\ge1$, the area of the region $(1+xy)/(x+y)\ge z$ is
\[
 \mathcal A(z)=\int_0^{1/z}\frac{1-zx}{z-x}\,dx
 =1-(z^2-1)\ln\frac{z^2}{z^2-1}
 =\sum_{k\ge1}\frac{z^{-2k}}{k(k+1)}.
\]
The endpoint $z=1$ is interpreted by a limit.
Layer cake and Tonelli's theorem give
\[
 I=\int_1^\infty\mathcal A(z)\,\frac{dz}{z}
   =\frac12\sum_{k\ge1}\frac1{k^2(k+1)}
   =\frac{\zeta(2)-1}{2}.
\]
The last expression uses
$1/(k^2(k+1))=1/k^2-1/k+1/(k+1)$ and Euler's classical
evaluation $\zeta(2)=\pi^2/6$.
\end{proof}

A primitive pair cannot have either $3$ or $5$ as a common divisor.
Since $f\ge0$, inclusion--exclusion in the larger set which excludes
only those two common divisors gives
\[
 P_f(T)\le F(T)-F(T/3)-F(T/5)+F(T/15).
\]
Dividing by an odd common divisor preserves the odd denominator.
Using the upper bound in \eqref{eq:grid} in the positive terms and
the lower bound in the negative terms yields
\[
 P_f(T)\le\frac{32I}{75}T^2+\frac8{15}T.
\]
This includes scaled radii smaller than $1$.
Substitution in \eqref{eq:U} gives the uniform finite estimate
\begin{equation}\label{eq:tail}
 \mathcal L_{\rm other}\le
 \frac{32I}{75}(1-2^{-J})+\frac{8J}{15q}.
\end{equation}

\section{Central credit and the infinite range}
For integer $T\ge1$, define the signed primitive count, with the axis once,
\[
 A(T)=\#\{(a,b):|a|\le T,\ 1\le b\le T,\ b\ \odd,\ \gcd(a,b)=1\}.
\]
\begin{lemma}\label{lem:count}
For every integer $T\ge1$,
\begin{equation}\label{eq:count}
 A(T)\ge\frac8{\pi^2}T^2-\frac34T\ln T-\frac{11}4T-\frac14.
\end{equation}
For $n\ge2$, $T_0=\lfloor q/\sqrt2\rfloor$ satisfies
$C_M(T_0)=A(T_0)$, and
\begin{equation}\label{eq:central}
 \mathcal D_M\ge\frac{2\ln2}{\pi^2}
 -\frac{4\sqrt2\ln2/\pi^2+
       (\ln2/(2\sqrt2))((3/4)\ln q+11/4)}q
 -\frac{\ln2}{8q^2}.
\end{equation}
\end{lemma}
\begin{proof}
M\"obius inversion over odd common divisors gives the exact formula
\begin{equation}\label{eq:mobius}
 A(T)=\sum_{\substack{g\le T\\g\ \odd}}
       \mu(g)(2\lfloor T/g\rfloor+1)
                   \left\lceil\frac{\lfloor T/g\rfloor}{2}\right\rceil.
\end{equation}
For $l=\lfloor x\rfloor$ the raw count
$B(l)=(2l+1)\lceil l/2\rceil$ equals $l^2+l/2$ when $l$ is even and
$l^2+3l/2+1/2$ when it is odd. Hence
$|B(l)-x^2|\le3x/2+1/2$.
Also
\[
 \sum_{\substack{g\le T\\g\ \odd}}\frac1g\le1+\frac12\ln T,
 \qquad \#\{g\le T:g\ \odd\}\le(T+1)/2.
\]
For the first inequality,
$1/(2k+1)\le\frac12\ln((2k+1)/(2k-1))$ for $k\ge1$ telescopes.
The raw error in \eqref{eq:mobius} is at most
$(3/4)T\ln T+(7/4)T+1/4$.
The absolutely convergent Euler product gives
$\sum_{g\ \odd}\mu(g)/g^2=8/\pi^2$; truncating it at $T$ costs at most
$T^2\sum_{g>T}g^{-2}\le T$.
This proves \eqref{eq:count}.

Two distinct primitive representations in the $T_0$ box of one residue
have a nonzero determinant divisible by $M$, unless they coincide.
The determinant has absolute value at most $2T_0B_0$, where $B_0$
is the largest odd integer at most $T_0$.
If $2T_0^2<M$, this is less than $M$.
If $2T_0^2=M$ and $n\ge2$, $T_0$ is a power of $2$ and is even, so
$B_0=T_0-1$ and the strict inequality still holds.
A zero determinant gives the same primitive fraction with positive
denominator. Thus the residue map is injective; every covered residue
has a primitive representation in the box, so $C_M(T_0)=A(T_0)$.
The excluded case $n=1$ has the collision $(1,1),(-1,1)$.

By monotonicity $\mathcal D_M\ge(\ln2)/(2M)A(T_0)$.
Use $T_0^2\ge q^2/2-\sqrt2q$,
$T_0\ln T_0\le(q/\sqrt2)\ln q$ and $T_0\le q/\sqrt2$ in
\eqref{eq:count}. This yields \eqref{eq:central}.
\end{proof}

The axis mass at valuation $k$ is $2^{-k-1}$. Consequently
\begin{equation}\label{eq:axis}
 \frac{\mathcal L_V}{\ln2}
 =\sum_{k=\lceil n/2\rceil}^{n-1}2^{-k-1}(k-n/2)
 =
 \begin{cases}
 2^{-\ell}-(\ell+1)2^{-2\ell},&n=2\ell,\ \ell\ge1,\\
 3\,2^{-\ell-2}-(2\ell+3)2^{-2\ell-2},&n=2\ell+1,\ \ell\ge0.
 \end{cases}
\end{equation}
In particular $\mathcal L_V\le\sqrt2\ln2/q$ for all $n\ge1$.
These are finite geometric sums, including the zero energies at $n=1,2$.

\begin{proposition}\label{prop:cutoff}
For every parent length $n\ge26$,
$\mathcal L_M<\mathcal D_M$, hence $K_M<\frac12\ln2$.
\end{proposition}
\begin{proof}
Put
\[
 c=\frac{32I}{75}=\frac{16}{75}\left(\frac{\pi^2}{6}-1\right),
 \qquad \delta=\frac{2\ln2}{\pi^2}-c.
\]
Equations \eqref{eq:tail}, \eqref{eq:central}, \eqref{eq:axis}, discarding
the favorable term $-c2^{-J}$, give
$\mathcal D_M-\mathcal L_M\ge\delta-E(q,J)$, where
\[
 E(q,J)=\frac{8J/15+4\sqrt2\ln2/\pi^2+
       (\ln2/(2\sqrt2))((3/4)\ln q+11/4)+\sqrt2\ln2}{q}
                   +\frac{\ln2}{8q^2}.
\]
Use $3<\pi<22/7$, $1<\sqrt2<3/2$, and
$\ln2>2/3+2/81=56/81$ from the positive atanh series. Then
\[
 \delta>\frac{1372}{9801}-\frac{304}{2205}
        =\frac{45756}{21611205}>\frac1{500}.
\]
Since $q\ge2^J$ and $\ln q<J+1$, the same elementary bounds give
\[
 E(q,J)\le
 \left(\frac{109}{120}J+\frac{47}{12}\right)2^{-J}
                   +\frac18\,2^{-2J}.
\]
At $J=13$ the right side is
$629/327680+1/536870912<1/520+1/13000=1/500$.
Both terms decrease for $J\ge13$; for the first,
$aJ+b>a$ with $a=109/120$, $b=47/12$.
Thus $n\ge26$ implies $E<1/500<\delta$.
\end{proof}

\section{Finite certificates and completion of the proof}\label{sec:finite}
The remaining premise is $K_{2^n}\le\frac12\ln2$ for $1\le n\le25$.
There are two different finite interfaces, described here so the computation
has an exact mathematical meaning. No decision uses floating logarithms.

For $n=1,\ldots,12$, certificates give the height $h_r$ and a primitive
witness $(a_r,b_r)$ for every residue. The checker verifies
$b_r>0$ odd, $a_r\equiv rb_r$, and
$\max(|a_r|,b_r)=h_r$. To prove minimality, for every positive odd $b<h_r$
it computes the least absolute representative of $rb\bmod M$ and
checks that it is at least $h_r$. This excludes every improving pair.
Let $c_n=\#\{r:2h_r^2>M\}$ and
$P_n=\prod_{2h_r^2>M}h_r^2$. Then the required comparison is exactly
\begin{equation}\label{eq:Kinteger}
                P_n\le 2^{\,M+(n-1)c_n}.
\end{equation}
This follows by exponentiating $2MK_M\le M\ln2$.
The full domain has $\sum_{n=1}^{12}2^n=8190$ minima.
For diagnostic support, the checker also verifies
$\prod_r h_r^2\le M^M$, but \eqref{eq:Kinteger} is the premise used
with the child lemma.

For $n=13,\ldots,25$, set $T_0=\lfloor\sqrt{M/2}\rfloor$ and compute
$A=A(T_0)$ exactly by \eqref{eq:mobius}. Put
\[
 V_2=\sum_{\substack{0\le k<n\\2k>n}}
                2^{n-k-1}(2k-n)=\frac{2M\mathcal L_V}{\ln2},
 \qquad
 x_{D,a,b}=\frac{(M+D^2ab)^2}{MD^2(a+b)^2}\ge1.
\]
Thus $2MU_M=\sum_{D,a,b}D\ln x_{D,a,b}$, with
$D=2^j$, $1\le j\le J$, $1\le a,b\le q/D$, $b$ odd, $\gcd(a,b)=1$.
Using $S=2^{40}$, let $\ell^+(x)\ge S\ln x$ and
$\ell^-_2\le S\ln2$ be the integer bounds of
Appendix~\ref{app:logs}. The sufficient certificate inequality is
\begin{equation}\label{eq:bridge}
 \sum_{D,a,b}D\ell^+(x_{D,a,b})\le(A-V_2)\ell^-_2,\qquad A>V_2.
\end{equation}
Indeed \eqref{eq:bridge} implies
$U_M+\mathcal L_V\le A\ln2/(2M)\le\mathcal D_M$.
The integer margin, right side minus left side, is listed below.
It is a scaled certificate margin, not the exact real difference
$\mathcal D_M-\mathcal L_M$.

\begin{center}
\begin{tabular}{rr@{\qquad}rr}
\toprule
Parent $n$&Integer margin&Parent $n$&Integer margin\\
\midrule
13&200913119387100&20&22717214856196764\\
14&343712368064520&21&45370620212212808\\
15&778422546045060&22&90744127844519936\\
16&1668697844541872&23&181329491713437076\\
17&2824453720758008&24&359749179298450700\\
18&5916038348384600&25&720718942077388912\\
19&11572174724235024&&\\
\bottomrule
\end{tabular}
\end{center}
The fixed panel contains $9\,068\,656$ positive primitive pairs, out of
$11\,182\,008$ raw pairs. The input and replay map in the companion separate
these essential certificates from later diagnostic controls.

\begin{proof}[Proof of Theorem~\ref{thm:main}]
The small certificates prove the sufficient bound on $K_M$ for
parents $n=1,\ldots,12$; the directed-log certificates prove it for
$n=13,\ldots,25$.
Proposition~\ref{prop:cutoff} proves it for $n\ge26$.
Equation~\eqref{eq:reduction} therefore gives
$\E\ln R_N\le(N/2)\ln2$ for every $N\ge2$, using parent length $n=N-1$.
For $N=1$, both residues have height $1$. Divide by $\ln2$.
\end{proof}

\section{Scope, reproducibility and remaining questions}
The theorem concerns the ordinary mean over all finite binary strings.
It gives no upper bound for every prescribed infinite $2$-adic number
and no statement about all Collatz trajectories.
The variance part of \cite[Question~3.1]{LXC}, sharper constants, and
extensions to other bases are not proved here.

The structural ingredient is the compensation between the two dyadic
lifts, followed by an upper bound in the actual odd completion class.
B\'ezout, Minkowski, M\"obius inversion, Euler products, convexity and
Tonelli are classical tools. The short count and high-valuation family
have explicit antecedents credited above.
The contribution is the exact mean bound and the estimates
which close its sign. Priority remains a separate source question.

The companion consists of this standalone source, a rendered PDF,
\path{companion/verify_mean.py}, \path{companion/finite_certificates.json}, a claim
map and reproduction instructions. From the repository root run
\begin{quote}\ttfamily
python -B scripts/reproduce.py\\
python -B -O scripts/reproduce.py
\end{quote}
The two runs must produce the same deterministic mathematical JSON.
The checker verifies level coverage, witnesses and minima, exact products,
directed-log sums, populations and margins, with explicit failures that
remain active under Python optimization.
It checks only the finite premises; reading the lemmas above is necessary
to assess the all-length proof.
The source compiles independently of the research reports and bibliography
databases.

\paragraph{AI assistance.}
OpenAI GPT models assisted mathematical exploration, proof criticism,
source comparison, editing and reproducibility checks.
The AI review passes shared exposure to the research material; they do
not constitute independent human peer review or formal verification.
The author retains responsibility for the arguments and attributions.

\appendix
\section{Directed logarithms using integer arithmetic}\label{app:logs}
We specify the rounding used in \eqref{eq:bridge}.
For a rational $v=u/w$ with $1\le v\le2$ let
$z=(u-w)/(u+w)\in[0,1/3]$. Then
\[
 \ln v=2\sum_{j=0}^{11}\frac{z^{2j+1}}{2j+1}
              +2\sum_{j=12}^\infty\frac{z^{2j+1}}{2j+1}.
\]
To compute an upper, replace $Sz$ by its ceiling, compute its square
divided by $S$ with upward rounding, start the scaled power with that
ceiling, and at each step round both division by $2j+1$ and the
multiplication by the scaled square divided by $S$ upwards.
Twice the twelve-term sum is an upper for the partial sum.
The exact tail, using $z\le1/2$, is bounded by
\[
 \frac{2z^{25}}{25(1-z^2)}
                  \le\frac8{3\cdot25\cdot2^{25}}.
\]
Add the integer ceiling of $8S/(3\cdot25\cdot2^{25})$.
For a lower, use downward rounding at every step and omit the
nonnegative remainder.
For $v=1$, return exactly zero.
These rules give a lower $\ell^-_2$ and upper $\ell^+_2$ for $S\ln2$.

For an arbitrary rational $x=u/w\ge1$, find the unique integer $k\ge0$
with $2^kw\le u<2^{k+1}w$, by integer bit lengths and a comparison.
Apply the unit-interval upper to $u/(2^kw)$ and add $k\ell^+_2$.
Every operation is an integer comparison, product, floor or ceiling.
Since all rounding maps are monotone on nonnegative inputs, the computed
scaled powers bound the true scaled powers by induction.
The uniform tail bounds the true analytic remainder; it does not assume
that an upward-rounded approximation of $z$ still equals $z$.
This proves the directed inequalities used by every summand of the
finite bridge.

\begin{thebibliography}{9}
\bibitem{TQ}
T. Tian and W.-F. Qi,
Expected values for the rational complexity of finite binary sequences,
\emph{Designs, Codes and Cryptography} \textbf{55} (2010), 65--79.
\href{https://doi.org/10.1007/s10623-009-9331-x}{doi:10.1007/s10623-009-9331-x}.
\bibitem{CWv1}
Z. Chen and A. Winterhof,
Probabilistic results on the 2-adic complexity,
arXiv:2501.16785v1, 28 January 2025.
\href{https://arxiv.org/abs/2501.16785v1}{arXiv version 1}.
Detailed theorem locators in this manuscript refer to this version.
\bibitem{CWfinal}
Z. Chen and A. Winterhof,
Probabilistic results on the 2-adic complexity,
\emph{Designs, Codes and Cryptography} \textbf{93} (2025), 2191--2203.
\href{https://doi.org/10.1007/s10623-025-01592-1}{doi:10.1007/s10623-025-01592-1}.
Version of record identified; detailed theorem locators above refer to arXiv v1.
\bibitem{LXC}
H. Lin, L. Xiao and Z. Chen,
A brief view on the linear and 2-adic complexities of pseudorandom binary sequences,
\emph{AIMS Mathematics} \textbf{11}(7) (2026), 20267--20290.
\href{https://doi.org/10.3934/math.2026823}{doi:10.3934/math.2026823}.
\bibitem{RW}
Y. Ren and A. Winterhof,
Symmetric measures of pseudorandomness for binary sequences,
arXiv:2603.23166v1, 24 March 2026.
\href{https://arxiv.org/abs/2603.23166v1}{arXiv version 1}.
\bibitem{GK}
M. Goresky and A. Klapper,
\emph{Algebraic Shift Register Sequences}, Cambridge University Press, 2012.
\href{https://doi.org/10.1017/CBO9781139057448}{doi:10.1017/CBO9781139057448}.
Growth lemma cited as Lemma~18.5.1 in \cite{CWv1}.
Public author draft, 14 October 2009, Section~21.5.a,
Lemma~21.5.1:
\href{https://www.cs.uky.edu/~klapper/pdf/algebraic.pdf}{author draft}.
Draft and published numbering are not identified here.
\bibitem{Cassels}
J. W. S. Cassels,
\emph{An Introduction to the Geometry of Numbers},
Springer, 1997 reprint, Chapter~III, Section~1 (Minkowski's convex body theorem).
\href{https://doi.org/10.1007/978-3-642-62035-5_4}{doi:10.1007/978-3-642-62035-5\_4}.
\end{thebibliography}
\end{document}

